% Abstract (NBR 6028: one paragraph, no formulas).
For a decade, model selection for tabular data relied on the maxim that
gradient boosting always wins. Tabular foundation models, pre-trained
networks that predict without task-specific training, have cast doubt on
that maxim, yet most of the available evidence comes from the models'
own proponents, under differing protocols. This dissertation aims to
produce comparative evidence under a single protocol and to turn it into
a selection procedure, along three lines. First, a multi-criteria
benchmark evaluates 14 models from three families (gradient boosting,
deep learning, and foundation models) on 18 public OpenML datasets, with
a hold-out that no selection stage ever observes, frequentist and
Bayesian inference, and measured training, tuning, and latency costs.
Second, a decision framework is assessed: its learned component by
leaving one dataset out at a
time (leave-one-dataset-out), its fixed policies descriptively on the
same datasets. Third, the distillation of the TabPFN
teacher's quantile curves into gradient-boosting students is studied on
20 regression datasets, under registered hypotheses. Among the 11
pre-2026 models, the top of the accuracy ranking is a tie: 3 to 8 of the
10 non-leading models are practically equivalent to the leader ---
itself already a foundation model, TabPFN ---,
depending on the Bayesian criterion adopted, within one percentage point
of AUC. In the 2026 generation, TabFM, with no tuning at all, has the
best average rank in all three tasks (in multiclass, over the two
datasets where it operates) and is the best model on 13 of the
18 datasets; the lead is consistent across datasets, but the tests,
underpowered with so few datasets per task, do not confirm it as a
statistical separation, and it costs four to five orders of magnitude
more inference latency. Kolmogorov--Arnold networks finish in the bottom
third on all three tasks. For the decision, a router learned from
meta-features is less accurate than the trivial rule (0.11 vs.\ 0.44);
the fixed ``foundation model first'' policy has median regret of zero on
the primary metric, but stops being the recommendation once 10 to 20
percent of the decision weight shifts from accuracy to latency (by
median regret), under the rank normalization adopted.
In distillation, the student trained on the teacher's means beats the
control on only 4 of 20 datasets; the student trained on the quantile
curves retains, at the median, 12 to 19 percent of the teacher's
distributional advantage, depending on how the students are tuned, with
positive retention on 16 of 20 datasets (sign test, p of approximately
0.019 under the two-stage design correction), and beats default-configured classical
distributional baselines on 17 of 20. The conclusion is that, in the
regime studied (datasets of up to 100 thousand samples, pinned model
versions, a single GPU), model choice should start from the
application's constraints (latency, architectural limits,
interpretability) rather than from a universal winner; and that
distillation recovers a modest median share of the foundation
model's predictive uncertainty, under preconditions that can be checked
in minutes.

\vspace{1cm}
\noindent\textbf{Keywords:} tabular data; foundation models;
benchmarking; knowledge distillation; distributional regression.
